% Authoritative exercise source. Preserve all derivations and diagrams.

\ExerciseStart{G2.06}{Leggere una matrice attraverso le immagini della base}{Consolidamento}

\BookHeading{Problema, dati e obiettivo}

Studiare la trasformazione \(A=\begin{pmatrix}1&2\\2&4\end{pmatrix}\). Calcolare le immagini della base canonica, del vettore \(v=(3,-1)^T\) e di \(w=(1,0)^T\); determinare nucleo e immagine.

\ExTheory{\ExRef{c2-matrici}{2.07}; \ExRef{c2-nucleo}{2.08}}{La colonna j di A è l’immagine di \(e_j\). Per \(x=(x_1,x_2)^T\), \(Ax=x_1Ae_1+x_2Ae_2\). Il nucleo contiene i vettori mandati in zero, l’immagine tutti i risultati possibili.}

\ExPrediction{Le colonne sono parallele. Il piano degli ingressi dovrebbe quindi essere mandato su una retta; una direzione intera perderà la propria componente distinguibile.}

\ExFigure{\begin{tikzpicture}[>=Stealth,exlabel]
\draw[exaxis] (-0.35,0)--(1.55,0) node[right] {$x_1$};\draw[exaxis] (0,-0.35)--(0,1.55) node[above] {$x_2$};\draw[->,Accent,line width=1.1pt] (0,0)--(1,0) node[anchor=south west] {$ e_1 $};\draw[->,Secondary,line width=1.1pt] (0,0)--(0,1) node[anchor=south west] {$ e_2 $};
\end{tikzpicture}}{Base del piano di ingresso.}


\begin{ExTasks}

\item Leggere e verificare le due colonne.

\item Eseguire il prodotto con v.

\item Eseguire il prodotto con w e confrontare gli ingressi.

\item Risolvere \(Ax=0\).

\item Parametrizzare l’immagine e spiegare il rango.

\end{ExTasks}

\ExSolution

\ExStep{Applicare A alla base}

\ExEq{Ae_1&=\colv{1\cdot1+2\cdot0\\2\cdot1+4\cdot0}=\colv{1\\2},\\Ae_2&=\colv{1\cdot0+2\cdot1\\2\cdot0+4\cdot1}=\colv{2\\4}=2Ae_1.}

\ExStep{Trasformare il primo ingresso}

\ExEq{Av&=\colv{1\cdot3+2(-1)\\2\cdot3+4(-1)}=\colv{3-2\\6-4}=\colv{1\\2}.}

\ExStep{Confrontare un altro ingresso}

\ExEq{Aw&=\colv{1\cdot1+2\cdot0\\2\cdot1+4\cdot0}=\colv{1\\2}.}
La differenza è \(v-w=(2,-1)^T\). Applicando A si ottiene \((2-2,4-4)^T=(0,0)^T\).

\ExStep{Risolvere il sistema del nucleo}

\ExEq{x_1+2x_2&=0,\qquad2x_1+4x_2=0.}
La seconda equazione è il doppio della prima. Posto \(x_2=t\), si ha \(x_1=-2t\), dunque \(\ker A=\{t(-2,1)^T:t\in\mathbb R\}\).

\ExStep{Ricavare l’immagine}

\ExEq{Ax&=\colv{x_1+2x_2\\2(x_1+2x_2)}=(x_1+2x_2)\colv{1\\2}.}
Ogni multiplo è ottenibile scegliendo \(x_2=0\). L’immagine è \(\operatorname{span}\{(1,2)^T\}\), di dimensione 1. La dimensione del nucleo è 1 e \(1+1=2\) verifica il teorema rango-nullità.

\par\noindent\begin{minipage}{\linewidth}
\ExFigure{\begin{tikzpicture}[>=Stealth,exlabel]
\draw[exaxis] (-0.35,0)--(2.55,0) node[right] {$x_1$};\draw[exaxis] (0,-0.35)--(0,4.55) node[above] {$x_2$};\draw[->,Accent,line width=1.1pt] (0,0)--(1,2) node[anchor=south west] {$ Ae_1 $};\draw[->,Secondary,line width=1.1pt] (0,0)--(2,4) node[anchor=south west] {$ Ae_2 $};
\end{tikzpicture}}{Le immagini delle due direzioni di base appartengono alla medesima retta.}


\ExInterpret{A conserva soltanto la combinazione \(x_1+2x_2\). Il determinante \(1\cdot4-2\cdot2=0\) conferma la perdita di invertibilità. Una trasformazione di questo tipo non rappresenta da sola l’evoluzione unitaria di un sistema quantistico chiuso.}

\ExLab{Nelle sezioni 2.07–2.08 inserire la matrice nello strumento delle trasformazioni e confrontare immagini della base, griglia trasformata e nucleo. Muovere l’ingresso lungo una retta parallela a \((-2,1)\): l’uscita resta invariata.}


\end{minipage}
